% Lecture 11 - Page 7 (13-Dec-1994).
% Fragment for % Lecture 11 - Page 7 (13-Dec-1994).
% Fragment for % Lecture 11 - Page 7 (13-Dec-1994).
% Fragment for % Lecture 11 - Page 7 (13-Dec-1994).
% Fragment for \input{courses/pages/curs11-pagina7-en}.
% Faithful translation of the manuscript; no preamble or forced page break.

\subsection*{Lecture 11 - Page 7}

\noindent\underline{Lemma:} Let $m,n,q\in\mathbb N^*$, $n=mq\Rightarrow n\mathbb Z\subseteq m\mathbb Z$.

\[
\frac{m\mathbb Z}{n\mathbb Z}\cong\mathbb Z_q.
\]

\noindent Define $f$ by the composition
\[
\begin{tikzcd}[column sep=large]
m\mathbb Z \arrow[r] \arrow[rr,bend left=30,"f"] &
\mathbb Z \arrow[r,"\mathrm{can}"] &
\mathbb Z/q\mathbb Z=\mathbb Z_q .
\end{tikzcd}
\]
\[
mk\longmapsto k\longmapsto \widehat{k}=k+q\mathbb Z.
\]
\[
m(h+k)\longmapsto h+k\longmapsto\widehat h+\widehat k.
\]

\noindent $f$ is a surjective homomorphism $\Rightarrow \operatorname{Im}f=\mathbb Z_q$.
\[
\mathbb Z_q\cong \frac{m\mathbb Z}{\ker f}.
\]

\noindent $x\in\ker f\Rightarrow f(x)=\widehat 0$.
\[
x\in m\mathbb Z,\quad x=mk\ \stackrel{f}{\longmapsto}\ \widehat k=\widehat 0
\iff q\mid k\iff k=qk_1.
\]
\[
\Rightarrow x=mqk_1=nk_1\quad\Rightarrow\quad\ker f=n\mathbb Z.
\]

\bigskip
\begin{center}
\fbox{\textbf{TABLE OF CONTENTS}}\qquad 2 theorems + 5 exercises
\end{center}
% The handwritten note appears to read "2 teorii + 5 exercitii".
% Number (5) is written twice / partly revised in the manuscript.
% The original numbering is preserved rather than silently regularised.

\begin{enumerate}
\item[(1)] Composition of mappings between sets. Properties.
\item[(2)] Equivalence relations, classes, transversals, quotient sets.
\item[(3)] Factorisation theorems for set maps.
\item[(4)] Internal laws of composition; associativity, commutativity, identity element.
\item[(5)] Composition law induced on a stable subset and on a quotient set.
\item[(5)] Monoid, group, subgroup: definitions, examples.
\item[(6)] Congruences on groups. Lagrange's theorem.
\item[(7)] Normal subgroup, quotient group.
\item[(8)] Order of an element, properties, examples.
\item[(9)] Group homomorphisms, examples, properties.
\item[(10)] Factorisation theorem for group homomorphisms, T.~F-I.
% "T. F-I" is the original abbreviation, probably referring to
% the fundamental isomorphism theorem; not expanded silently.
\item[(11)] The symmetric group $S_n$, transpositions, cycles; decomposition of a
permutation into a product of transpositions (cycles).
% Item (11) continues at the beginning of Page 8.
\end{enumerate}
.
% Faithful translation of the manuscript; no preamble or forced page break.

\subsection*{Lecture 11 - Page 7}

\noindent\underline{Lemma:} Let $m,n,q\in\mathbb N^*$, $n=mq\Rightarrow n\mathbb Z\subseteq m\mathbb Z$.

\[
\frac{m\mathbb Z}{n\mathbb Z}\cong\mathbb Z_q.
\]

\noindent Define $f$ by the composition
\[
\begin{tikzcd}[column sep=large]
m\mathbb Z \arrow[r] \arrow[rr,bend left=30,"f"] &
\mathbb Z \arrow[r,"\mathrm{can}"] &
\mathbb Z/q\mathbb Z=\mathbb Z_q .
\end{tikzcd}
\]
\[
mk\longmapsto k\longmapsto \widehat{k}=k+q\mathbb Z.
\]
\[
m(h+k)\longmapsto h+k\longmapsto\widehat h+\widehat k.
\]

\noindent $f$ is a surjective homomorphism $\Rightarrow \operatorname{Im}f=\mathbb Z_q$.
\[
\mathbb Z_q\cong \frac{m\mathbb Z}{\ker f}.
\]

\noindent $x\in\ker f\Rightarrow f(x)=\widehat 0$.
\[
x\in m\mathbb Z,\quad x=mk\ \stackrel{f}{\longmapsto}\ \widehat k=\widehat 0
\iff q\mid k\iff k=qk_1.
\]
\[
\Rightarrow x=mqk_1=nk_1\quad\Rightarrow\quad\ker f=n\mathbb Z.
\]

\bigskip
\begin{center}
\fbox{\textbf{TABLE OF CONTENTS}}\qquad 2 theorems + 5 exercises
\end{center}
% The handwritten note appears to read "2 teorii + 5 exercitii".
% Number (5) is written twice / partly revised in the manuscript.
% The original numbering is preserved rather than silently regularised.

\begin{enumerate}
\item[(1)] Composition of mappings between sets. Properties.
\item[(2)] Equivalence relations, classes, transversals, quotient sets.
\item[(3)] Factorisation theorems for set maps.
\item[(4)] Internal laws of composition; associativity, commutativity, identity element.
\item[(5)] Composition law induced on a stable subset and on a quotient set.
\item[(5)] Monoid, group, subgroup: definitions, examples.
\item[(6)] Congruences on groups. Lagrange's theorem.
\item[(7)] Normal subgroup, quotient group.
\item[(8)] Order of an element, properties, examples.
\item[(9)] Group homomorphisms, examples, properties.
\item[(10)] Factorisation theorem for group homomorphisms, T.~F-I.
% "T. F-I" is the original abbreviation, probably referring to
% the fundamental isomorphism theorem; not expanded silently.
\item[(11)] The symmetric group $S_n$, transpositions, cycles; decomposition of a
permutation into a product of transpositions (cycles).
% Item (11) continues at the beginning of Page 8.
\end{enumerate}
.
% Faithful translation of the manuscript; no preamble or forced page break.

\subsection*{Lecture 11 - Page 7}

\noindent\underline{Lemma:} Let $m,n,q\in\mathbb N^*$, $n=mq\Rightarrow n\mathbb Z\subseteq m\mathbb Z$.

\[
\frac{m\mathbb Z}{n\mathbb Z}\cong\mathbb Z_q.
\]

\noindent Define $f$ by the composition
\[
\begin{tikzcd}[column sep=large]
m\mathbb Z \arrow[r] \arrow[rr,bend left=30,"f"] &
\mathbb Z \arrow[r,"\mathrm{can}"] &
\mathbb Z/q\mathbb Z=\mathbb Z_q .
\end{tikzcd}
\]
\[
mk\longmapsto k\longmapsto \widehat{k}=k+q\mathbb Z.
\]
\[
m(h+k)\longmapsto h+k\longmapsto\widehat h+\widehat k.
\]

\noindent $f$ is a surjective homomorphism $\Rightarrow \operatorname{Im}f=\mathbb Z_q$.
\[
\mathbb Z_q\cong \frac{m\mathbb Z}{\ker f}.
\]

\noindent $x\in\ker f\Rightarrow f(x)=\widehat 0$.
\[
x\in m\mathbb Z,\quad x=mk\ \stackrel{f}{\longmapsto}\ \widehat k=\widehat 0
\iff q\mid k\iff k=qk_1.
\]
\[
\Rightarrow x=mqk_1=nk_1\quad\Rightarrow\quad\ker f=n\mathbb Z.
\]

\bigskip
\begin{center}
\fbox{\textbf{TABLE OF CONTENTS}}\qquad 2 theorems + 5 exercises
\end{center}
% The handwritten note appears to read "2 teorii + 5 exercitii".
% Number (5) is written twice / partly revised in the manuscript.
% The original numbering is preserved rather than silently regularised.

\begin{enumerate}
\item[(1)] Composition of mappings between sets. Properties.
\item[(2)] Equivalence relations, classes, transversals, quotient sets.
\item[(3)] Factorisation theorems for set maps.
\item[(4)] Internal laws of composition; associativity, commutativity, identity element.
\item[(5)] Composition law induced on a stable subset and on a quotient set.
\item[(5)] Monoid, group, subgroup: definitions, examples.
\item[(6)] Congruences on groups. Lagrange's theorem.
\item[(7)] Normal subgroup, quotient group.
\item[(8)] Order of an element, properties, examples.
\item[(9)] Group homomorphisms, examples, properties.
\item[(10)] Factorisation theorem for group homomorphisms, T.~F-I.
% "T. F-I" is the original abbreviation, probably referring to
% the fundamental isomorphism theorem; not expanded silently.
\item[(11)] The symmetric group $S_n$, transpositions, cycles; decomposition of a
permutation into a product of transpositions (cycles).
% Item (11) continues at the beginning of Page 8.
\end{enumerate}
.
% Faithful translation of the manuscript; no preamble or forced page break.

\subsection*{Lecture 11 - Page 7}

\noindent\underline{Lemma:} Let $m,n,q\in\mathbb N^*$, $n=mq\Rightarrow n\mathbb Z\subseteq m\mathbb Z$.

\[
\frac{m\mathbb Z}{n\mathbb Z}\cong\mathbb Z_q.
\]

\noindent Define $f$ by the composition
\[
\begin{tikzcd}[column sep=large]
m\mathbb Z \arrow[r] \arrow[rr,bend left=30,"f"] &
\mathbb Z \arrow[r,"\mathrm{can}"] &
\mathbb Z/q\mathbb Z=\mathbb Z_q .
\end{tikzcd}
\]
\[
mk\longmapsto k\longmapsto \widehat{k}=k+q\mathbb Z.
\]
\[
m(h+k)\longmapsto h+k\longmapsto\widehat h+\widehat k.
\]

\noindent $f$ is a surjective homomorphism $\Rightarrow \operatorname{Im}f=\mathbb Z_q$.
\[
\mathbb Z_q\cong \frac{m\mathbb Z}{\ker f}.
\]

\noindent $x\in\ker f\Rightarrow f(x)=\widehat 0$.
\[
x\in m\mathbb Z,\quad x=mk\ \stackrel{f}{\longmapsto}\ \widehat k=\widehat 0
\iff q\mid k\iff k=qk_1.
\]
\[
\Rightarrow x=mqk_1=nk_1\quad\Rightarrow\quad\ker f=n\mathbb Z.
\]

\bigskip
\begin{center}
\fbox{\textbf{TABLE OF CONTENTS}}\qquad 2 theorems + 5 exercises
\end{center}
% The handwritten note appears to read "2 teorii + 5 exercitii".
% Number (5) is written twice / partly revised in the manuscript.
% The original numbering is preserved rather than silently regularised.

\begin{enumerate}
\item[(1)] Composition of mappings between sets. Properties.
\item[(2)] Equivalence relations, classes, transversals, quotient sets.
\item[(3)] Factorisation theorems for set maps.
\item[(4)] Internal laws of composition; associativity, commutativity, identity element.
\item[(5)] Composition law induced on a stable subset and on a quotient set.
\item[(5)] Monoid, group, subgroup: definitions, examples.
\item[(6)] Congruences on groups. Lagrange's theorem.
\item[(7)] Normal subgroup, quotient group.
\item[(8)] Order of an element, properties, examples.
\item[(9)] Group homomorphisms, examples, properties.
\item[(10)] Factorisation theorem for group homomorphisms, T.~F-I.
% "T. F-I" is the original abbreviation, probably referring to
% the fundamental isomorphism theorem; not expanded silently.
\item[(11)] The symmetric group $S_n$, transpositions, cycles; decomposition of a
permutation into a product of transpositions (cycles).
% Item (11) continues at the beginning of Page 8.
\end{enumerate}
